\documentclass[11pt]{article}

\usepackage[a4paper,margin=1in]{geometry}
\usepackage{array}
\usepackage{booktabs}
\usepackage{enumitem}
\usepackage[T1]{fontenc}
\usepackage{hyperref}
\usepackage{longtable}

\newcommand{\ecfile}[1]{\texttt{#1}}
\newcommand{\ecname}[1]{\texttt{\detokenize{#1}}}
\newcommand{\claimstatus}[1]{\textbf{#1}}

\title{EasyCrypt Proof-Script Commentary for the HAETAE Provable-Security Artifact}
\author{}
\date{May 2026}

\begin{document}
\maketitle

\section{Purpose}

This note records manuscript-ready commentary for the EasyCrypt proof
scripts in the HAETAE provable-security artifact.  The goal is to state
precisely what the checked development establishes, what remains
conditional, and which proof obligations should be targeted next.

\paragraph{Verification command.}
The checked proof surface is the manifest
\ecfile{provable-security/proof-files.txt}.  The repository verification
gate is:
\begin{quote}
\ecname{sh provable-security/verify-provable-security.sh}
\end{quote}
The gate compiles each listed EasyCrypt file with include paths
\ecfile{provable-security/easycrypt} and \ecfile{kyber-security}.  In the
current run, all sixteen manifest files compiled successfully and the
script ended with \ecname{RESULT: PASS}.

\section{Honest Machine-Checked Claim}

The current EasyCrypt development machine-checks a formal HAETAE model,
perfect correctness for that model, and a family of conditional EUF-CMA
reduction theorems in the random-oracle setting.  The strongest checked
EUF-CMA theorem bounds the adversary success probability by the counted
HAETAE ROM bound, assuming an explicit exact-hyperball NMA bound and
abstract MLWE and Module-SIS hardness bounds.  The artifact does not prove
unconditional cryptographic security, does not discharge the cryptographic
hardness assumptions, and does not contain a HAETAE SUF-CMA theorem beyond
the generic SUF-CMA game definition.

\section{Claim Matrix}

\begin{longtable}{>{\raggedright\arraybackslash}p{0.30\linewidth}
                  >{\raggedright\arraybackslash}p{0.36\linewidth}
                  >{\raggedright\arraybackslash}p{0.24\linewidth}}
\toprule
Claim & EasyCrypt item & Status \\
\midrule
\endhead
HAETAE signing and verification are perfectly correct in the EasyCrypt
model.
&
\ecname{correctness_perfect},
\ecfile{HAETAE\_Security.ec:40}
&
\claimstatus{Machine-checked, model-conditional.}
\\
\addlinespace
The EUF-CMA experiment is defined.
&
\ecname{module EUF_CMA},
\ecfile{Sig\_ROM.ec:94}
&
\claimstatus{Game/interface definition only.}
\\
\addlinespace
The SUF-CMA experiment is defined.
&
\ecname{module SUF_CMA},
\ecfile{Sig\_ROM.ec:174}
&
\claimstatus{Definition only.}  No HAETAE SUF-CMA theorem was found in the
checked proof surface.
\\
\addlinespace
The abstract HAETAE EUF-CMA bound follows from a Fiat--Shamir hop, an NMA
bound, a bimodal-to-MSIS hop, and MLWE/Module-SIS hardness bounds.
&
\ecname{euf_cma_security},
\ecfile{HAETAE\_Security.ec:136}
&
\claimstatus{Conditional/abstract reduction.}
\\
\addlinespace
The counted-ROM HAETAE EUF-CMA bound follows from the exact-hyperball NMA
premise and abstract hardness premises.
&
\ecname{euf_cma_security_from_checked_ro_sampling_hop_and_counted_bimodal},
\ecfile{HAETAE\_Security.ec:1217}
&
\claimstatus{Machine-checked conditional reduction.}
\\
\addlinespace
The budget-explicit path exposes the signing-budget adapters and the
sampling-loss obligation.
&
\ecname{euf_cma_security_from_budgeted_paper_sample_lifting_and_counted_bimodal},
\ecfile{HAETAE\_Security.ec:1170}
&
\claimstatus{Conditional; exposes proof obligations.}
\\
\addlinespace
The bimodal self-target MSIS game can be related to Module-SIS.
&
\ecname{bimodal_to_module_sis_reduction_bound},
\ecfile{HAETAE\_Security.ec:418}
&
\claimstatus{Checked wrapper, hardness-conditional.}
\\
\bottomrule
\end{longtable}

\section{Dependency Chain for the Strongest EUF-CMA Theorem}

The strongest theorem in the current checked surface is
\ecname{euf_cma_security_from_checked_ro_sampling_hop_and_counted_bimodal}
in \ecfile{HAETAE\_Security.ec:1217}.  Its conclusion is:
\begin{quote}
\ecname{Pr[SIG.EUF_CMA(H, HAETAE, A).main() @ &m : res] <= haetae_euf_counted_rom_bound}.
\end{quote}
The theorem depends on the following chain.

\begin{enumerate}[leftmargin=*]
\item The theorem applies
\ecname{euf_cma_security_from_ro_exact_hyperball_sampling_hop_and_counted_bimodal}
(\ecfile{HAETAE\_Security.ec:1036}).

\item The real-signing to RO exact-hyperball hop is supplied by
\ecname{real_signing_ro_exact_hyperball_hop_checked_from_loss_budget}
(\ecfile{HAETAE\_Security.ec:1102}).

\item That lemma applies
\ecname{real_signing_ro_exact_hyperball_hop_from_ro_signing_attempt_hop}
(\ecfile{HAETAE\_Security.ec:865}) and then calls
\ecname{rom_internal_nma_ro_signing_attempt_ro_exact_hyperball_loss_bound}
(\ecfile{HAETAE\_HopGames.ec:7554}).

\item The current proof of
\ecname{rom_internal_nma_ro_signing_attempt_ro_exact_hyperball_loss_bound}
is a coarse fallback: it uses probability boundedness and
\ecname{rejection_sampling_loss_term_ge1} rather than an adversary-level
sampling hybrid.  This makes the statement checked, but not a non-vacuous
cryptographic sampling-loss proof.

\item The final EUF-CMA theorem also assumes the exact-hyperball NMA
premise:
\begin{quote}
\ecname{Pr[UF_NMA(... ROExactHyperballPaperSimSampler ...)] + rejection_sampling_loss_term <= Pr[MLWE_Game(Bmlwe)] + Pr[BimodalSelfTargetMSIS_Game(Bbimodal)]}.
\end{quote}
This premise is not discharged by the theorem itself.

\item The MLWE and Module-SIS terms are abstract hardness premises:
\ecname{mlwe_hardness_term} and \ecname{module_sis_hardness_term} in
\ecfile{HAETAE\_Reductions.ec:9--10}.
\end{enumerate}

\section{Trust Boundaries}

\paragraph{Abstract types and distributions.}
The model declares \ecname{mlwe_instance} as abstract and introduces
lossless distributions such as \ecname{dmlwe_real},
\ecname{dmlwe_random}, \ecname{dmodule_sis_instance}, and
\ecname{dbimodal_selftarget_msis_instance}
(\ecfile{HAETAE\_Assumptions.ec:9--19}).  These are part of the formal
model boundary.

\paragraph{Hardness terms.}
The quantities \ecname{mlwe_hardness_term} and
\ecname{module_sis_hardness_term} are abstract real-valued bounds
(\ecfile{HAETAE\_Reductions.ec:9--10}).  EasyCrypt checks reductions to
these terms; it does not establish the concrete hardness of MLWE or
Module-SIS.

\paragraph{Declared modules.}
The theorem sections declare adversaries and reductions, for example
\ecname{H}, \ecname{A}, \ecname{Bmlwe}, \ecname{Bbimodal}, and
\ecname{Bsis} in \ecfile{HAETAE\_Security.ec:129--134}.  The checked
statements are parametric over these interfaces.

\paragraph{Coarse bounds.}
Several lemmas use the fact that a probability is at most one together
with \ecname{rejection_sampling_loss_term_ge1}.  The important remaining
coarse sites are
\ecfile{HAETAE\_HopGames.ec:7305},
\ecfile{HAETAE\_HopGames.ec:7402},
\ecfile{HAETAE\_HopGames.ec:7554}, and
\ecfile{HAETAE\_HopGames.ec:9593}.  These are checked EasyCrypt facts,
but they should be described as fallback or vacuous bounds, not as tight
cryptographic reductions.

\section{Recent Proof-Script Improvement}

The latest proof improvement extends the bridge chain in
\ecfile{provable-security/easycrypt/HAETAE\_HopGames.ec}.  The first is
\ecname{concrete_ro_signing_attempt_to_exact_hyperball_sample_with_seed_clean_state_loss}
at line 6844.  It lifts the existing clean one-call sampler loss to
predicates that may mention the final random-oracle global state, provided
the proof supplies left and right framing inequalities.  The second is
\ecname{concrete_budgeted_o_sign_sample_with_seed_clean_state_loss_surface}
at line 6912, specialized to the budgeted NMA adapter's ROM query logs.
The latest addition is
\ecname{concrete_budgeted_o_sign_clean_stateful_loss_from_sample_frames}
at line 8034.  It composes the state-aware sampler bridge with explicit
attempt-side and exact-side \ecname{O.sign} frame inequalities, proving that
one clean active signing call can be lifted with one
\ecname{rejection_sampling_loss_term} once the continuation has already been
summarized as predicates over the returned signature, the sampled
\ecname{paper_sim_signature_sample}, and the final random-oracle state.

These lemmas are non-vacuous because they reuse the checked one-call loss
instead of relying on \ecname{rejection_sampling_loss_term_ge1}.  They are
still bridge or surface lemmas: they do not yet replace the adversary-level
NMA fallback at \ecfile{HAETAE\_HopGames.ec:7554}, and the concrete clean
NMA lifting at \ecfile{HAETAE\_HopGames.ec:9593} remains coarse.

\paragraph{Current first-active two-call work.}
The proof script now contains a seed-sensitive suffix attempt for the
actual first-active two-call \ecname{O.sign} branch.  The new EasyCrypt
surface introduces \ecname{concrete\_fro\_get\_step}, a pure model of the
checked lazy random-oracle \ecname{FRO.get} map update, and
\ecname{concrete\_two\_call\_first\_active\_o\_sign\_suffix\_kernel},
which records the sampler-expand query, reconstructs the sampler ROM
entry from the sampled signing coins, runs the message-hash and
challenge-hash suffix, computes the signature/transcript, and evaluates
the full-state postcondition.  The seeded bridge
\ecname{structural\_to\_exact\_hyperball\_seeded\_coin\_sample\_pair\_suffix\_loss\_from\_point\_loss}
keeps the stronger point-loss premise and does not use the coarse NMA
target.  Verification is pending in this edit pass; the commentary should
only be upgraded to a checked claim after \ecname{HAETAE\_HopGames.ec} and
the full manifest compile successfully.

\section{Publication Wording}

\paragraph{Acceptable wording.}
We machine-checked in EasyCrypt a HAETAE random-oracle model, a perfect
correctness theorem, and conditional EUF-CMA reduction chains with explicit
MLWE and Module-SIS hardness assumptions and stated sampling/ROM trust
boundaries.

\paragraph{Too-strong wording to avoid.}
The manuscript should not say that HAETAE is fully machine-checked secure,
unconditionally secure, SUF-CMA verified, or that every cryptographic
reduction and sampling hop has been discharged non-vacuously.

\paragraph{Precise statement.}
The EasyCrypt artifact verifies a conditional EUF-CMA theorem for the
modeled HAETAE signature scheme: assuming the stated exact-hyperball NMA
bound and abstract MLWE/Module-SIS hardness bounds, the adversary's
EUF-CMA success probability is bounded by the counted HAETAE ROM security
expression.  The proof is machine-checked over the listed EasyCrypt files,
but the final cryptographic claim remains conditional on the abstract
hardness assumptions, declared reduction interfaces, and several explicit
sampling/ROM boundaries, including coarse fallback bounds that should not
be presented as tight reductions.

\section{Recommended Next Codex Request}

The highest-value next request is to target the smallest adversary-level
gap that the new bridge lemmas make approachable.  The request should also
make the manuscript file a live work product: every proof-status change must
be reflected in this document before the agent stops.

\begin{quote}
\small
You are a Codex AI agent in the current HAETAE EasyCrypt project.  Task:
replace the coarse EasyCrypt proof of the budgeted clean NMA sampling
lifting with a non-vacuous proof using the checked one-call sampling-loss
bridges.  Start from
\ecname{concrete_rom_internal_budgeted_nma_ro_signing_attempt_ro_exact_hyperball_clean_conditioned_lifting}
in \ecfile{provable-security/easycrypt/HAETAE\_HopGames.ec}.  Do not make
broad refactors.  Preserve existing theorem statements unless a minimal
strengthening is necessary.  Use
\ecname{concrete_budgeted_o_sign_sample_with_seed_clean_state_loss_surface}
and related projection lemmas to prove the clean conditioned NMA bound by
lifting one \ecname{O.sign} call and accumulating at most
\ecname{rom_signature_query_budget * rejection_sampling_loss_term}.  Avoid
using \ecname{rejection_sampling_loss_term_ge1} or probability-boundedness
as the main proof.

Continuously maintain the manuscript commentary while working.  Update
\ecfile{manuscript/haetae-easycrypt-proof-script-commentary.tex} whenever a
proof status changes: before editing, record the target and current trust
boundary; after each successful proof step, update the relevant claim matrix,
dependency chain, and recent proof-script improvement text; after any failed
attempt, record the blocker only if it changes the proof plan or trust
boundary; after verification, update the commands/results and publication
wording if the honest claim changed.  Keep the LaTeX concise and
manuscript-ready, not as a raw work log.

Rerun \ecname{sh provable-security/verify-provable-security.sh} and compile
the manuscript LaTeX if \ecname{pdflatex} is available.  Report exactly which
coarse-bound dependency was removed, which premises remain, which manuscript
sections were updated, and any new proof obligations.
\end{quote}

\section{Continuous Documentation Protocol}

Proof development and manuscript maintenance should proceed together.  The
document is not merely an after-the-fact summary: it is the running record of
the current honest security claim.  A Codex agent working on the proof should
update the manuscript at four checkpoints:

\begin{enumerate}[leftmargin=*]
\item \textbf{Before proof edits:} identify the target lemma, the current
fallback or premise it relies on, and the expected non-vacuous replacement.

\item \textbf{After a checked proof step:} update the claim matrix and
dependency chain if a lemma changes status from coarse, premise-level, or
abstract to non-vacuously proved.

\item \textbf{After a failed proof attempt:} update the proof plan only when
the failure reveals a real missing lemma, interface mismatch, or trust
boundary.  Do not turn transient tactic failures into manuscript content.

\item \textbf{After verification:} record the exact EasyCrypt verification
command, the result, the changed theorem status, and any change to acceptable
publication wording.
\end{enumerate}

\subsection*{Current Proof-Edit Checkpoint}

The active proof target is
\ecname{concrete_rom_internal_budgeted_nma_ro_signing_attempt_ro_exact_hyperball_clean_conditioned_lifting}
in \ecfile{provable-security/easycrypt/HAETAE\_HopGames.ec}.  Its current
proof is a coarse fallback: it bounds the clean conditioned NMA event by
one and then uses
\ecname{rom_signature_query_budget * rejection_sampling_loss_term >= 1}.
The intended replacement is a non-vacuous lifting argument that applies the
checked one-call bridge
\ecname{concrete_budgeted_o_sign_sample_with_seed_clean_state_loss_surface}
inside the budgeted signing oracle and accumulates at most
\ecname{rom_signature_query_budget * rejection_sampling_loss_term} over the
checked signing-query budget.

\subsection*{Proof Attempt Status}

A new checked local support lemma,
\ecname{concrete_budgeted_o_sign_clean_stateful_loss_from_sample_frames},
now proves the one-call stateful frame step.  Its proof is non-vacuous: it
uses
\ecname{concrete_budgeted_o_sign_sample_with_seed_clean_state_loss_surface}
and does not appeal to \ecname{rejection_sampling_loss_term_ge1} or
probability boundedness.

The direct replacement of the target NMA lemma is still blocked by an
adversary-continuation gap rather than by a local sampling lemma.  The new
stateful frame lemma handles one active \ecname{O.sign} call after the
pre-call state and post-call continuation have already been summarized by
explicit frame inequalities.  The target NMA lemma needs a stronger rule:
after an active signing call, the postcondition is the rest of the adaptive
adversary execution plus the final UF-NMA verifier.  That continuation
depends on \ecname{glob A}, \ecname{glob HAETAE_RO.FRO}, the
transcript/signing logs, the query counters, and the clean flag, not only on
the returned signature.

Consequently, replacing the target proof non-vacuously requires one of the
following checked ingredients before the target lemma can be closed:
\begin{enumerate}[leftmargin=*]
\item a stateful approximate call rule for
\ecname{ROMInternalTranscriptBudgetedPaperSimAsNMA(...).O.sign} that accepts
continuation predicates over the full post-call state and adds one
\ecname{rejection_sampling_loss_term}; or
\item a concrete hybrid wrapper, analogous to EasyCrypt's hybrid theorem,
that packages adversary hash queries as leaks, packages signing as the
hybridized oracle, proves the one-switch stateful signing lemma, and sums
the loss over \ecname{signature_query_budget_count} active signing calls.
\end{enumerate}

The current target lemma therefore remains in the trust boundary until that
adversary-level stateful lifting principle is proved.  Replacing its body
with a proof that merely routes through another coarse lemma, or again uses
\ecname{rejection_sampling_loss_term_ge1}, would not change the honest
security claim.

The current follow-up isolates and checks the saturated half of that missing
interface.  The proof script now contains
\ecname{concrete_budgeted_o_sign_saturated_fallback_equiv}
(\ecfile{HAETAE\_HopGames.ec:8672}),
\ecname{concrete_budgeted_o_sign_saturated_fallback_full_state_equiv}
(\ecfile{HAETAE\_HopGames.ec:8701}),
\ecname{concrete_budgeted_ah_get_full_state_equiv}
(\ecfile{HAETAE\_HopGames.ec:8758}),
\ecname{concrete_budgeted_adversary_forge_saturated_fallback_equiv}
(\ecfile{HAETAE\_HopGames.ec:8832}), and
\ecname{concrete_budgeted_adversary_forge_saturated_clean_fallback_lifts}
(\ecfile{HAETAE\_HopGames.ec:8920}).  Together these lemmas show that,
once \ecname{signature_query_budget_count <= signing_count}, the adaptive
\ecname{A.forge} continuation can keep making hash and signing calls without
incurring additional sampler loss: hash calls are simulated identically, and
post-budget signing calls take the same fallback branch in the signing-attempt
and exact-hyperball worlds while preserving the ROM, transcript, query-log,
and clean-flag state.  The new probability-lifting wrapper converts the
saturated equivalence into the exact inequality needed for predicates of the
form \ecname{R res (glob HAETAE_RO.FRO)} under the clean event.

This is a checked adversary-level saturated-continuation fact, but it is not
yet the active-call accumulation theorem.  The remaining interface is the
one-switch active-call hybrid: before saturation, choose one active call
satisfying \ecname{signing_count < signature_query_budget_count}, apply
\ecname{concrete_budgeted_o_sign_clean_stateful_loss_from_sample_frames}
(\ecfile{HAETAE\_HopGames.ec:8034}) to that call, carry the full ROM and
transcript state through the remaining \ecname{A.forge} continuation, then sum
exactly one \ecname{rejection_sampling_loss_term} for each active call.  No
unchecked one-switch lemma was kept in the source.

The current active-call follow-up adds an explicit checked two-call hybrid
wrapper for the concrete route:
\ecname{ROMInternalTranscriptBudgetedTwoCallHybridAsNMA}
(\ecfile{HAETAE\_HopGames.ec:1766}), plus the three intended instances
\ecname{ROMInternalTranscriptBudgetedTwoCallG0AsNMA}
(\ecfile{HAETAE\_HopGames.ec:1879}),
\ecname{ROMInternalTranscriptBudgetedTwoCallG1AsNMA}
(\ecfile{HAETAE\_HopGames.ec:1885}), and
\ecname{ROMInternalTranscriptBudgetedTwoCallG2AsNMA}
(\ecfile{HAETAE\_HopGames.ec:1891}).  The wrapper keeps the same ROM,
transcript, query-log, and clean-flag state as the budgeted game, but selects
\ecname{First.sample_with_seed} on the first active call and
\ecname{Second.sample_with_seed} on the second active call.  Thus
\ecname{G0} uses signing-attempt sampling twice, \ecname{G1} switches the
first active call to exact-hyperball sampling, and \ecname{G2} uses
exact-hyperball sampling twice.  These are now real EasyCrypt games, not just
comments or planned notation.

The same follow-up retains two checked counter facts for the concrete
two-call route:
\ecname{budgeted_paper_sim_sign_oracle_first_active_sets_one}
(\ecfile{HAETAE\_HopGames.ec:1926}) and
\ecname{budgeted_paper_sim_sign_oracle_second_active_saturates}
(\ecfile{HAETAE\_HopGames.ec:1943}).  These lemmas use
\ecname{signature_query_budget_count = 2} to show that the first active
\ecname{O.sign} call moves \ecname{signing_count} from \ecname{0} to
\ecname{1}, and the second active call moves it from \ecname{1} to the
saturated budget value.  They justify reducing the intended hybrid proof to
two active switches followed by the already checked saturated continuation,
but they do not by themselves compose the sampling loss through the adaptive
continuation.

The current first-switch follow-up proves the exact continuation needed after
the first active signing call, but does not yet prove the approximate
\ecname{G0 -> G1} bound itself.  Once the shared two-call state satisfies
\ecname{1 <= signing_count}, \ecname{G0} and \ecname{G1} no longer use the
different first-call sampler: the next active call uses the signing-attempt
\ecname{Second} sampler on both sides, and post-budget calls use the same
fallback branch.  This is checked at the hash-oracle level by
\ecname{concrete_two_call_g0_g1_post_first_ah_get_equiv}
(\ecfile{HAETAE\_HopGames.ec:9002}), at the signing-oracle level by
\ecname{concrete_two_call_g0_g1_post_first_o_sign_equiv}
(\ecfile{HAETAE\_HopGames.ec:9080}), and through the adaptive adversary
continuation by
\ecname{concrete_two_call_g0_g1_post_first_adversary_forge_equiv}
(\ecfile{HAETAE\_HopGames.ec:9171}).  The probability wrapper
\ecname{concrete_two_call_g0_g1_post_first_clean_adversary_lifts}
(\ecfile{HAETAE\_HopGames.ec:9264}) packages this as a clean-event
continuation fact with zero additional sampler loss.  The remaining missing
principle was the approximate split at the first active \ecname{O.sign} call:
apply the one-call rejection-sampling loss to that call and then resume with
the checked post-first continuation.

The latest first-active-call follow-up adds the checked frame-composition
lemma
\ecname{concrete_two_call_g0_g1_first_active_o_sign_clean_stateful_loss_from_sample_frames}
(\ecfile{HAETAE\_HopGames.ec:8509}).  This is the two-call-wrapper analogue
of the one-call stateful bridge: under clean ROM coverage, equal sampler
secret-key state, and \ecname{signing_count = 0}, it turns explicit
attempt-side and exact-side sample-frame premises into a one
\ecname{rejection_sampling_loss_term} bound from
\ecname{ROMInternalTranscriptBudgetedTwoCallG0AsNMA.O.sign} to
\ecname{ROMInternalTranscriptBudgetedTwoCallG1AsNMA.O.sign}.  Its predicate
surface includes the returned signature, the final ROM global, and the full
two-call adapter state
(\ecname{pk_current}, \ecname{queries}, \ecname{transcripts},
\ecname{records}, \ecname{adversary_hash_count}, \ecname{signing_count},
\ecname{adversary_hash_queries}, \ecname{sampler_expand_queries}, and
\ecname{sampler_bad_prequery}), and the event records
\ecname{1 <= signing_count}.  This makes the lemma continuation-ready, but
the general full-state sample-frame inequalities remain premises of this
generic lemma.

The newest concrete frame step discharges the signature-level first-active
\ecname{O.sign} switch without using \ecname{mu_bounded},
\ecname{rejection_sampling_loss_term_ge1}, or the coarse target lemma.  The
attempt side is checked by
\ecname{concrete_two_call_g0_first_active_o_sign_active_clean_signature_structural_le}
(\ecfile{HAETAE\_HopGames.ec:8236}) and
\ecname{concrete_two_call_g0_first_active_o_sign_to_self_logged_sample_active_projection}
(\ecfile{HAETAE\_HopGames.ec:8386}).  The exact side is checked by
\ecname{concrete_two_call_g1_first_active_o_sign_active_signature_exact}
(\ecfile{HAETAE\_HopGames.ec:8319}) and
\ecname{concrete_two_call_g1_exact_sample_with_seed_to_first_active_o_sign_projection}
(\ecfile{HAETAE\_HopGames.ec:8466}).  These projections are then packaged in
\ecname{concrete_two_call_g0_g1_first_active_o_sign_clean_signature_loss}
(\ecfile{HAETAE\_HopGames.ec:8614}), which gives a concrete
\ecname{G0.O.sign <= G1.O.sign + rejection_sampling_loss_term} bound for
signature predicates under \ecname{signing_count = 0}, clean ROM coverage,
and equal sampler secret-key state.  The remaining adversary-level gap is to
lift that first-call switch to predicates over the full post-call
ROM/transcript/two-call adapter state and then apply
\ecname{concrete_two_call_g0_g1_post_first_clean_adversary_lifts} as the
adaptive continuation.

The latest full-state-lift pass checked the seed-sensitive bridge that was
missing from the signature-only switch.  The helper
\ecname{concrete_mu_dlet_ge_point} (\ecfile{HAETAE\_HopGames.ec:8672}) proves
that one point contribution is bounded by the probability of the whole
continuation bind.  Using that helper,
\ecname{structural_to_exact_hyperball_coin_sample_pair_suffix_loss_from_point_loss}
(\ecfile{HAETAE\_HopGames.ec:8683}) lifts the rejection-sampling point-loss
obligation through an arbitrary common randomized suffix over
\ecname{signing_attempt_coin_sample_pair}.  This bridge is seed-sensitive:
the suffix may observe the sampled signing seed as well as the sampled paper
simulation value.

The checked wrapper
\ecname{concrete_two_call_g0_g1_first_active_o_sign_clean_full_state_loss_from_coin_sample_suffix_frames}
(\ecfile{HAETAE\_HopGames.ec:8743}) packages that bridge for the desired
full-state \ecname{G0.O.sign <= G1.O.sign + rejection_sampling_loss_term}
surface.  Its final predicate ranges over the returned signature,
\ecname{glob HAETAE\_RO.FRO}, \ecname{pk_current}, query/transcript/record
logs, counters, adversary and sampler query logs, and
\ecname{sampler_bad_prequery}.  The lemma is still a frame-composition result:
it requires concrete left and right frame inequalities connecting the actual
two-call \ecname{O.sign} code to a common suffix kernel \ecname{K}.  Those
frame inequalities are now the exact remaining obligation for the full-state
first-call switch.

This new bridge deliberately uses the stronger
\ecname{structural_to_exact_hyperball_sample_pair_point_loss_obligation},
not only the paper-sample predicate obligation.  The stronger premise is
needed because the full post-call state can depend on the seed recorded in
\ecname{sampler_expand_queries}, and because the two post-sampler
\ecname{FRO.get} calls make the final ROM global a randomized continuation of
the coin/sample pair rather than a boolean predicate on
\ecname{paper_sim_signature_sample}.  No new unchecked EasyCrypt assumption
was added.  The concrete full-state \ecname{O.sign} switch, the adversary-level
\ecname{G0 -> G1} lift, and the coarse NMA target lemma remain unchanged.

The exact endpoint of the two-call hybrid is now connected to the existing
exact-hyperball budgeted game by checked equivalences.  The proof script first
aligns the hash oracle in
\ecname{concrete_two_call_g2_exact_budgeted_ah_get_equiv}
(\ecfile{HAETAE\_HopGames.ec:9312}), then aligns the signing oracle in
\ecname{concrete_two_call_g2_exact_budgeted_o_sign_equiv}
(\ecfile{HAETAE\_HopGames.ec:9389}).  The signing-oracle proof explicitly
case-splits on the active/saturated branch and on whether the active call is
the first or second one, forcing the wrapper's
\ecname{First}/\ecname{Second} dispatch before simulation.  These two oracle
equivalences are lifted through \ecname{A.forge} by
\ecname{concrete_two_call_g2_exact_budgeted_adversary_forge_equiv}
(\ecfile{HAETAE\_HopGames.ec:9462}), through the wrapper's public
\ecname{forge} procedure by
\ecname{concrete_two_call_g2_exact_budgeted_forge_equiv}
(\ecfile{HAETAE\_HopGames.ec:9554}), and finally through the NMA game by
\ecname{concrete_two_call_g2_exact_budgeted_nma_equiv}
(\ecfile{HAETAE\_HopGames.ec:9569}).  This closes the exact \ecname{G2}
endpoint; the remaining work is the adversary-level \ecname{G0 -> G1}
lift, now starting from a checked signature-level first-call switch, and the
second active-call switch \ecname{G1 -> G2}.

\subsection*{Verification Checkpoint}

After the current follow-up, the EasyCrypt source retains the checked
stateful one-call frame lemma, the checked first-active two-call
sample-frame composition lemma, the checked concrete first-active
signature-level \ecname{G0}/\ecname{G1} \ecname{O.sign} loss lemma, the
checked explicit two-call hybrid games, the checked post-first
\ecname{G0}/\ecname{G1} continuation, the checked exact
\ecname{G2} endpoint equivalence, the
checked active counter lemmas, the checked saturated adversary-continuation
equivalence, and the checked saturated clean probability-lifting wrapper.  The
edited EasyCrypt file was locally checked
with:
\begin{quote}
\ecname{easycrypt compile provable-security/easycrypt/HAETAE_HopGames.ec -max-provers 1 -I provable-security/easycrypt -I kyber-security}
\end{quote}
The file compiled successfully.  The full proof gate was rerun with:
\begin{quote}
\ecname{sh provable-security/verify-provable-security.sh}
\end{quote}
The command compiled all sixteen manifest files successfully and ended with
\ecname{RESULT: PASS}.  The target clean NMA lifting theorem
\ecname{concrete_rom_internal_budgeted_nma_ro_signing_attempt_ro_exact_hyperball_clean_conditioned_lifting}
(\ecfile{HAETAE\_HopGames.ec:9593}) remains machine-checked only as a coarse
fallback; the checked one-call stateful frame lemma, explicit two-call hybrid
games, first-active two-call sample-frame composition lemma, concrete
first-active signature-level \ecname{O.sign} switch, post-first
\ecname{G0}/\ecname{G1} continuation, exact \ecname{G2} endpoint
equivalence, active counter lemmas, and saturated adversary-continuation
lemmas are non-vacuous pieces, but they have not yet removed that fallback.
This manuscript file was also recompiled with
\ecname{pdflatex -interaction=nonstopmode -halt-on-error}; compilation
succeeded, with only overfull-box warnings from long EasyCrypt identifiers.

\subsection*{Seed-Sensitive Suffix-Frame Follow-Up}

The current seed-sensitive suffix-frame follow-up has not yet produced a
checked full-state first-active \ecname{O.sign} switch.  The active proof
attempt instantiates the seeded suffix-kernel bridge shape for the actual
two-call first-active branch.  The \ecname{G0} frame
\ecname{concrete\_two\_call\_g0\_first\_active\_o\_sign\_clean\_full\_state\_frame\_to\_seeded\_coin\_sample\_suffix}
now compiles past its previous rejected post-sampler pointwise
\ecname{mu\_eq} blocker.  The local repair adds a checked pointwise
post-sampler two-\ecname{FRO.get} helper,
\ecname{rejected\_clean\_signature\_event\_two\_fro\_get\_drop\_count\_suffix\_false},
which matches the generated count-dropped clean rejection predicate without
using coarse \ecname{mu\_bounded} or
\ecname{rejection\_sampling\_loss\_term\_ge1} reasoning.

The most recent full compile command was:
\begin{quote}
\ecname{easycrypt compile provable-security/easycrypt/HAETAE\_HopGames.ec -max-provers 1 -I provable-security/easycrypt -I kyber-security}
\end{quote}
This run passed the \ecname{G0} frame and then failed with
\ecname{cannot prove goal (strict)} in the next downstream
\ecname{G1} frame lemma,
\ecname{concrete\_two\_call\_g1\_first\_active\_o\_sign\_clean\_full\_state\_frame\_from\_seeded\_coin\_sample\_suffix}.
Therefore the full-state first-active \ecname{G0 -> G1}
\ecname{O.sign} switch and the adversary-level lift through
\ecname{concrete\_two\_call\_g0\_g1\_post\_first\_clean\_adversary\_lifts}
remain unchecked.  The stronger
\ecname{structural\_to\_exact\_hyperball\_sample\_pair\_point\_loss\_obligation}
premise is still present, and the coarse NMA target lemma has not been
changed.

\section{Prioritized Proof Plan}

\begin{enumerate}[leftmargin=*]
\item Replace
\ecname{concrete_rom_internal_budgeted_nma_ro_signing_attempt_ro_exact_hyperball_clean_conditioned_lifting}
(\ecfile{HAETAE\_HopGames.ec:9593}).  Expected edits:
\ecfile{provable-security/easycrypt/HAETAE\_HopGames.ec}.  Difficulty:
high.  Strategy: prove an adversary-call hybrid over the budgeted
\ecname{O.sign} oracle, apply
\ecname{concrete_budgeted_o_sign_clean_stateful_loss_from_sample_frames} at
the selected clean active signing call, use
\ecname{budgeted_paper_sim_sign_oracle_first_active_sets_one} and
\ecname{budgeted_paper_sim_sign_oracle_second_active_saturates} to split the
two active switches, use
\ecname{concrete_budgeted_adversary_forge_saturated_clean_fallback_lifts} for
the post-budget continuation, and accumulate the loss by the checked
signing-count budget.

\item Replace
\ecname{rom_internal_nma_ro_signing_attempt_ro_exact_hyperball_loss_bound}
(\ecfile{HAETAE\_HopGames.ec:7554}).  Expected edits:
\ecfile{provable-security/easycrypt/HAETAE\_HopGames.ec} and possibly
\ecfile{HAETAE\_Security.ec} if the final theorem should use the
budget-explicit path.  Difficulty: medium-high.  Strategy: route through
budgeted clean lifting and the bad-prequery finite-event-loss bound rather
than through \ecname{rejection_sampling_loss_term_ge1}.

\item Discharge the exact-hyperball NMA-to-hardness premise in
\ecname{euf_cma_security_from_checked_ro_sampling_hop_and_counted_bimodal}
(\ecfile{HAETAE\_Security.ec:1217}).  Expected edits:
\ecfile{HAETAE\_HopGames.ec}, \ecfile{HAETAE\_Transcript.ec},
\ecfile{HAETAE\_Assumptions.ec}, and \ecfile{HAETAE\_Security.ec}.
Difficulty: high.  Strategy: connect the exact-hyperball NMA game to the
MLWE and bimodal self-target MSIS games via the existing transcript,
special-soundness, and bimodal-to-Module-SIS machinery.
\end{enumerate}

\end{document}
